\chapter{Gauss equations and curvature inequalities}\label{sec:gauss-simons}
Let $\Sigma^n$ be a smooth two-sided immersed hypersurface of a space form of sectional curvature $c$, with $n\geq2$. We use $h(X,Y)=\langle\overline\nabla_X\nu,Y\rangle$, $H=\operatorname{tr}h$, and $A=h^\sharp$. Thus a Euclidean sphere with outward normal has positive principal curvatures. The Laplacian is $\Delta=\operatorname{tr}\nabla^2$. These conventions agree with \cite{HP99}, Section 2, pp.~47--49. All calculations below are local; completeness, embeddedness, and compactness are unnecessary until an integral conclusion explicitly requires them.

\begin{proposition}[Gauss and Codazzi in a space form]\label{prop:spaceform-gauss-codazzi}
With $R(X,Y)Z=\nabla_X\nabla_YZ-\nabla_Y\nabla_XZ-\nabla_{[X,Y]}Z$, one has
\[
 R(X,Y)Z=c\bigl(\langle Y,Z\rangle X-\langle X,Z\rangle Y\bigr)
       +h(Y,Z)AX-h(X,Z)AY,
\]
\[
 (\nabla_Xh)(Y,Z)=(\nabla_Yh)(X,Z).
\]
In a principal orthonormal frame, the sectional curvature of the plane spanned by $e_i,e_j$, $i\ne j$, is $K_{ij}=c+\kappa_i\kappa_j$. Compare \cite{HP99}, Section 2, p.~49.
\end{proposition}
\begin{proof}
Differentiating $\langle\nu,Y\rangle=0$ gives
$\overline\nabla_XY=\nabla_XY-h(X,Y)\nu$, while $\overline\nabla_X\nu=AX$. Insert these decompositions into the ambient curvature expression. Its tangential part is
\[
 (\overline R(X,Y)Z)^T=R(X,Y)Z-h(Y,Z)AX+h(X,Z)AY.
\]
The ambient space-form formula yields the stated Gauss equation. The normal part is
$-(\nabla_Xh)(Y,Z)+(\nabla_Yh)(X,Z)$ and vanishes because the ambient curvature expression is tangent. This is Codazzi. Finally take $X=e_i$, $Y=Z=e_j$, and pair with $e_i$ in the Gauss equation.
\end{proof}

\begin{lemma}[Commuting derivatives of a covariant two-tensor]\label{lem:two-tensor-commutation}
For a smooth covariant two-tensor $T$, define $(\nabla^2T)(X,Y)=\nabla_X\nabla_YT-\nabla_{\nabla_XY}T$. Then
\[
 \bigl((\nabla^2T)(X,Y)-(\nabla^2T)(Y,X)\bigr)(Z,W)
 =-T(R(X,Y)Z,W)-T(Z,R(X,Y)W).
\]
\end{lemma}
\begin{proof}
Apply $XY-YX-[X,Y]$ to the scalar function $T(Z,W)$; the result is zero. Expand each derivative by the Leibniz rule for the connection. Terms with one derivative on each of two factors cancel in pairs. The terms with two derivatives on $Z$ and $W$ are respectively $T(R(X,Y)Z,W)$ and $T(Z,R(X,Y)W)$. The remaining terms give the derivative commutator on $T$. Moving the two curvature terms to the other side proves the signs in the formula. Torsion-freeness gives $\nabla_XY-\nabla_YX=[X,Y]$.
\end{proof}

\begin{theorem}[Simons identity in a space form]\label{thm:spaceform-simons}
For the hypersurface above, without any assumption that $H$ is constant,
\[
 \Delta h_{ij}=\nabla_i\nabla_jH+Hh_i{}^kh_{kj}-|A|^2h_{ij}
                 +c(nh_{ij}-Hg_{ij}).
\]
This is the space-form specialization of \cite{HP99}, Corollary 2.2, p.~49.
\end{theorem}
\begin{proof}
At the point under consideration choose a geodesic orthonormal frame that diagonalizes $h$, with eigenvalues $\kappa_i$. Codazzi and its contraction $\sum_k\nabla_kh_{kj}=\nabla_jH$ give
\[
 (\Delta h)_{ij}=\nabla_i\nabla_jH
 +\sum_k\bigl((\nabla^2h)(e_k,e_i)-(\nabla^2h)(e_i,e_k)\bigr)(e_k,e_j).
\]
The Hessian here includes the connection correction in Lemma \ref{lem:two-tensor-commutation}; it is not an uncorrected second derivative in a moving frame. By that lemma the remaining sum is
\[
 -\sum_k h(R(e_k,e_i)e_k,e_j)-\sum_kh(e_k,R(e_k,e_i)e_j).
\]
The Gauss equation makes this zero for $i\ne j$. For $i=j$ it equals
\[
 \sum_{k\ne i}(c+\kappa_k\kappa_i)(\kappa_i-\kappa_k)
 =c(n\kappa_i-H)+H\kappa_i^2-|A|^2\kappa_i.
\]
Indeed, the $c$ terms sum to $c(n\kappa_i-H)$, and the other terms use $\sum_k\kappa_k=H$ and $\sum_k\kappa_k^2=|A|^2$; the omitted $k=i$ term is zero. This proves every component in a diagonalizing basis and hence the tensor identity in all coordinates.
\end{proof}

\begin{corollary}[The norm identity with variable mean curvature]\label{cor:simons-norm-variable}
The same hypersurface satisfies
\[
 \frac12\Delta|A|^2=|\nabla A|^2+\langle h,\nabla^2H\rangle
 +H\operatorname{tr}(A^3)-|A|^4+c(n|A|^2-H^2).
\]
In particular, for a Euclidean minimal hypersurface,
$\frac12\Delta|A|^2=|\nabla A|^2-|A|^4$.
\end{corollary}
\begin{proof}
Metric compatibility gives $\frac12\Delta|h|^2=|\nabla h|^2+\langle h,\Delta h\rangle$. Contract Theorem \ref{thm:spaceform-simons} with $h$. The contractions with $h^2$ and $g$ are $\operatorname{tr}(A^3)$ and $H$, respectively. Setting $c=H=0$ proves the specialization.
\end{proof}

\begin{lemma}[Kato inequality and zeros of a tensor]\label{lem:tensor-kato}
For a smooth tensor $T$ on a Riemannian manifold, $u=|T|$ is locally Lipschitz and has a weak gradient satisfying $|\nabla u|\leq|\nabla T|$ almost everywhere. On $\{u>0\}$,
$\partial_i u=\langle\nabla_iT,T\rangle/u$.
\end{lemma}
\begin{proof}
Parallel transport along short curves is an isometry; the difference of the norms of the endpoint values of $T$ is bounded by the integral of $|\nabla T|$. Local boundedness of $\nabla T$ proves local Lipschitz regularity. Theorem \ref{thm:rademacher}, in coordinate charts, supplies almost-everywhere differentiability and the weak gradient. At a differentiability point where $u=0$, the nonnegative function $u$ has a local minimum and its differential is zero. At the other points differentiate $u^2=\langle T,T\rangle$ and use Cauchy--Schwarz, then sum over an orthonormal frame. This also justifies the inequality on the zero set, where division by $u$ is unavailable.
\end{proof}

\begin{lemma}[Improved Kato inequality for constant mean curvature]\label{lem:cmc-improved-kato}
If $H$ is constant on a hypersurface in a space form, then, wherever $|A|>0$,
\[
 |\nabla A|^2\geq\left(1+\frac2n\right)|\nabla|A||^2.
\]
The same assertion holds with $A$ replaced by its trace-free part $A-(H/n)I$ where that tensor is nonzero. Compare the Euclidean minimal case of \cite{SSY75}, equations (1.29)--(1.33), pp.~279--280; the constant-mean-curvature extension follows from the same algebra below.
\end{lemma}
\begin{proof}
At a point diagonalize the tensor under consideration, with eigenvalues $\lambda_i$. Its derivative $T_{ijk}=\nabla_kh_{ij}$ is fully symmetric by symmetry of $h$ and Codazzi, and $\sum_iT_{iik}=\partial_kH=0$. For fixed $k$, write $a_i=T_{iik}$. Then $a_k=-\sum_{i\ne k}a_i$, so Cauchy--Schwarz implies
\[
 a_k^2\leq\frac{n-1}{n}\sum_i a_i^2.
\]
Each component with exactly two equal indices occurs three times in $|T|^2$, and each component with three equal indices occurs once. Discarding the nonnegative contributions from three distinct indices therefore gives
\[
 |T|^2\geq\sum_k\left(3\sum_iT_{iik}^2-2T_{kkk}^2\right)
 \geq\left(1+\frac2n\right)\sum_{i,k}T_{iik}^2.
\]
On the other hand $\partial_k|h|=\sum_i\lambda_iT_{iik}/|h|$, so $|\nabla|h||^2\leq\sum_{i,k}T_{iik}^2$. This proves the first claim. When $H$ is constant, subtracting $(H/n)g$ leaves the derivative unchanged and gives another Codazzi tensor with constant trace; the same argument proves the second claim.
\end{proof}

\begin{proposition}[Simons inequality away from flat points]\label{prop:minimal-simons-inequality}
For a smooth Euclidean minimal hypersurface of dimension $n\geq2$, put $u=|A|$. On $\{u>0\}$,
\[
 u\Delta u+u^4\geq\frac2n|\nabla u|^2.
\]
No stability assumption is required for this differential inequality.
\end{proposition}
\begin{proof}
The identity $\frac12\Delta u^2=u\Delta u+|\nabla u|^2$, Corollary \ref{cor:simons-norm-variable}, and Lemma \ref{lem:cmc-improved-kato} give
$u\Delta u+u^4=|\nabla A|^2-|\nabla u|^2\geq(2/n)|\nabla u|^2$. Restriction to $u>0$ ensures the classical differentiability used here.
\end{proof}

\begin{proposition}[A weak curvature inequality through flat points]\label{prop:weak-curvature-subsolution}
On every smooth Euclidean minimal hypersurface, $u=|A|$ satisfies $\Delta u\geq-u^3$ in the distributional sense: for each nonnegative $\phi\in C_c^\infty(\Sigma)$,
\[
 -\int_\Sigma\langle\nabla u,\nabla\phi\rangle\geq-\int_\Sigma u^3\phi.
\]
The support of the test function lies in the smooth interior if the hypersurface has a boundary.
\end{proposition}
\begin{proof}
Put $u_\varepsilon=(|A|^2+\varepsilon^2)^{1/2}$. The norm identity gives
\[
 \Delta u_\varepsilon
 =\frac{|\nabla A|^2-u^4}{u_\varepsilon}
   -\frac{|\nabla(u^2)|^2}{4u_\varepsilon^3}
 \geq-\frac{u^4}{u_\varepsilon}\geq-u^3,
\]
because $|\nabla(u^2)|^2\leq4u^2|\nabla A|^2$. Multiply by $\phi$, integrate, and use $\int u_\varepsilon\Delta\phi=\int\phi\Delta u_\varepsilon$. Uniform convergence on the compact support gives $\int u\Delta\phi\geq-\int u^3\phi$. Lemma \ref{lem:tensor-kato} identifies the first term with the weak-gradient expression. No boundary term appears.
\end{proof}

\begin{proposition}[A curvature Caccioppoli estimate]\label{prop:curvature-caccioppoli}
Let $\Sigma^n$ be a smooth Euclidean minimal hypersurface, let $p\geq2$, and let $\phi\in C_c^\infty(\Sigma)$ have support in its interior. With $u=|A|$,
\[
 \int\phi^2u^{p-2}|\nabla A|^2
 \leq2\int\phi^2u^{p+2}+4\int u^p|\nabla\phi|^2,
\]
\[
 \int\phi^2|\nabla(u^{p/2})|^2
 \leq\frac{p^2}{2}\int\phi^2u^{p+2}+p^2\int u^p|\nabla\phi|^2.
\]
For $p=2$ the factor $u^{p-2}$ is interpreted as one, including at zeros. These estimates alone do not give an epsilon-regularity theorem.
\end{proposition}
\begin{proof}
Test $\Delta h=-u^2h$ against $\phi^2u_\varepsilon^{p-2}h$ and integrate by parts. Differentiating the weight yields a nonnegative term
$(p-2)\phi^2u_\varepsilon^{p-4}\sum_k\langle h,\nabla_kh\rangle^2$ on the gradient side. Dropping it gives
\[
 \int\phi^2u_\varepsilon^{p-2}|\nabla h|^2
 \leq\int\phi^2u_\varepsilon^{p-2}u^4
 +2\int |\phi|u_\varepsilon^{p-2}u|\nabla h||\nabla\phi|.
\]
Young's inequality bounds the last integrand by half the left gradient integrand plus $2u_\varepsilon^{p-2}u^2|\nabla\phi|^2$. Absorb the half and let $\varepsilon\downarrow0$. All functions are bounded on the fixed compact support and $p\geq2$, so dominated convergence proves the first estimate. The weak chain rule and Lemma \ref{lem:tensor-kato} give $|\nabla(u^{p/2})|^2\leq(p^2/4)u^{p-2}|\nabla A|^2$ almost everywhere, including the zero set. Multiplication by $p^2/4$ proves the second estimate.
\end{proof}
